%! TeX root = ./main.tex
% vim: set spelllang=en:
\chapter*{Foreword}
\addcontentsline{toc}{chapter}{Foreword}
\markboth{\textsc{foreword}}{}

% Phylogénétique, introduite par l’exemple de l’évolution des gènes
% Exemples de motivations
Genes are the carriers of heredity in living organisms.
They contribute to defining their traits and can be passed on from one generation to the next.
Over time and through various events, they can evolve into different alleles and spread across multiple species.
Phylogenetics aims to trace back the evolutionary history of units such as genes, cells, or species.
These histories are valuable sources of information that can help elucidate gene functions~\parencite{Lynch2001ProbabilityPreservationNewly}, characterize the evolution of tumours~\parencite{Quinn2021SingleCellLineages} or that of pathogens~\parencite{Morel2020PhylogeneticAnalysisSars,Guindon2022EclairerOriginesCOVID19}.

% Arbres phylogénétiques et événements
A \emph{phylogenetic tree} (or \emph{phylogeny}) represents a hypothesis of evolutionary history.
Much like a family tree, a phylogenetic tree is made up of a succession of branching points and of branches along which time flows until reaching the contemporary representatives.
Each branching point of the tree corresponds to the putative divergence of an ancestor into several descendants that are free to evolve independently of one another.
The possible causes of such divergences differ according to the types of units under study; for example, genes duplicate, cells divide and species undergo speciations.

% Problème d’évolution hiérarchique
The evolution of a gene depends not only on the events directly affecting it, but also on the evolution of its host species.
Indeed, the speciation of a species causes the emergence of independent copies of its genes, one in each new species.
This relationship forms an evolutionary hierarchy between a set of \emph{evolutionary units} and their hosts.
Similar hierarchies can be found at other scales, for example between protein domains and their proteins, parasites and their hosts, or species and their habitats.
Such hierarchies also arise outside of biology, for example, between cultural elements and the populations that transmit them \parencite{Tehrani2010CophylogenyPopulationsCultures}.

% Réconciliation
To study these evolutionary hierarchies, \textcite{Goodman1979FittingGeneLineage} introduced the so-called \emph{reconciliation} method.
Relying on the topological differences and similarities between the phylogeny of a gene family and that of its host species, reconciliation makes it possible to establish which of the two each divergence relates to.
The result is a reconciled tree, a synthesis of the two trees, in which each branching point of the gene tree is annotated with the type of evolutionary event that gave rise to it.

% Problème d’évolution segmentale
When a host carries multiple families of units, their phylogenies can be reconciled individually.
This assumes, however, that the units in question have evolved independently of one another.
This assumption is invalidated in particular when the units are organized linearly within the same host, such as genes along a genome, and are affected by large-scale events.
One of the research directions of the field over the past two decades has been the integration of reconciliation methods with those allowing the study of the evolution of linear structures~\parencite{Sankoff2000DuplicationRearrangementReconciliation}.
Various extensions of the reconciliation paradigm that account for the joint evolution of several units have thus been proposed.

% Projet de recherche
The research project carried out in this thesis falls within this research direction, through the development of \emph{multiple reconciliation} methods for studying the evolution of highly conserved genome segments called \emph{syntenies}.
The contribution notably consists in identifying trade-offs between modelling power and algorithmic complexity.

\Cref{chap:reconciliation} contains a review of the relevant literature that synthesizes existing reconciliation methods.
In it, we introduce the segmental reconciliation approach that lies at the heart of the research presented thereafter.
In \cref{chap:superdtl}, we make a first contribution by allowing the inference of horizontal transfers in segmental reconciliation.
This contribution is extended to the context of unsampled species in \cref{chap:synesth}.
These extensions lead us to the need for an in-depth combinatorial study of the solution spaces associated with the problems under study.
To achieve this, we develop a unified theory of algebraic dynamic programming in \cref{chap:adp}.
This theory is put into practice on the segmental reconciliation problem in \cref{chap:adp-rec}.
Finally, the various software tools and programs developed in support of the research conducted in this thesis are presented in \cref{chap:logiciel}.
